% Auto-generated by scripts/06_emit_structure.py
% Source: scans 020–044
% Section id: chapter/1
% Phase 3 deliverable. Footnotes (Phase 4), citations (Phase 5),
% and Wade-Giles → Pinyin (Phase 6) are not yet applied here.

\chapter{\pinyinterm{shang} Divination Procedures}
\label{ch:1}

% source: scan 020, printed 3
% intro-echo-dropped: \pinyinterm{shang}
% intro-echo-dropped: Divination Procedures
Pyro-Scapulimancy as a Culture Trait
Divining the future by the colorations, cracks, and other features of animal
shoulder blades is a human custom widespread and venerable.\footnote[1]{On scapulimancy (also called omoplatoscopy) in general, Eisenberger (1938) is the basic source; Andrée (1906) is still useful. When they touch on the custom, the articles in Caquot and Leibovici (1968) cite the more modern scholarship.}¹ Historical accounts often
fail to distinguish between apyro-scapulimancy, on the one hand, in which the natural
condition of the bone was read after the flesh had been scraped away, and pyro-scapu-limancy, on the other, in which the diviner created omen cracks in the bone by roasting
or burning. As a result, the exact extent of the culture area covered by the interventionist,
pyro-scapulimancy-that which the \pinyinterm{shang} practiced-is not clear. Generally speaking,
however, the scapulimancy of Europe, the Near East, and North Africa was apyromantic,
that of North and Central Asia and North America, pyromantic.\footnote[2]{Cooper (1936), pp. 34, 38. Eisenberger (1938), pp. 105, 111-112, notes only two instances of European pyromancy, in Hungary and in Scotland. Lin Sheng (1963) records the presence of both forms of scapulimancy in South China. On the genetic relation between the two forms, see Cooper (1936), pp. 39-40, whose terms ``pyro-`` and ``apyro-scapulimancy'' I borrow here; Eisenberger, disagreeing with Cooper, believes apyroscapulimancy to have been the older form.}²
Pyro-scapulimancy, narrowly construed to include only those cases in which a diviner
read the stress cracks caused by catastrophic heating of shoulder blades, was first practiced
by the Neolithic inhabitants of North China on the bones of deer, sheep, pigs, and cattle.
The earliest evidence for its existence, in Liaoning, may be dated toward the last half of
the fourth millennium B.C. By the end of the next millennium, scapulimancy had become
one of the characteristic features of Lungshanoid village life, the charred scapulimantic
remains suggesting the appearance of specialized groups of diviners and of a pyromantic
theology associated, in all probability, with ancestor worship.\footnote[3]{Itō (1962a), pp. 251-252, and \pinyinterm{zhang-bingquan} (1967), pp. 842-851, list pyromantic finds from 龙 and Bronze Age sites; cf. \pinyinterm{liu-yuanlin} (1974). For scapulimancy, Lungshanoid culture, and ancestor worship, see \pinyinterm{zhang-guangzhi} (1962), p. 1843 (1968), p. 128; Wheatley (1971), pp. 27-28. For carbon-14 dates of Lung3 ] ] ] 1 1 ] ] 1 7 ] 0}³
% source: scan 021, printed 4
There is impressive evidence for the custom's proliferation in historical times. To the
\pinyinterm{shang} kings, as we shall see, pyromantic reassurance was the sine qua non of daily life,
and consulting the turtle continued to play a significant, if diminished, role in China
during imperial times. * Close to home, Chinese emissaries reported the custom from Yamato
Japan in the third century A.D.,³ and it is mentioned in Kojiki and Manyōshū passages of
the eighth century.\footnote[4]{The history and social significance of pyromancy in China, to say nothing of that of Chinese divination in general, has yet to be written. For Eastern \pinyinterm{zhou-dynasty} plastromancy, see Yagimoto (1966); for the \pinyinterm{han-dynasty} situation, see 史赤, ch. 128. Hayashi Taisuke (1909), pp. 1243-1244, lists the plastromantic works cited in the bibliographical sections of later dynastic histories. Fujino (1960), p. 22, n. 6, lists the various plastromantic officers recorded down to 唐 times; as he notes (pp. 3-4), plastromancy appears to have been a technical speciality, centered on the court, but probably also practiced by the common people.}\footnote[5]{Tsunoda (1951), p. 12. For Yayoi scapulimancy, see Saitō (1958), pp. 242-243.}\footnote[6]{Chamberlain (1883), p. lix; Satow (1879), pp. 429-433; Philippi (1968), p. 52, n. 1.}\footnote[7]{For the reasons why Jordanes' record that Attila's haruspices ``examined the entrails of cattle and certain streaks in the bones that had been scraped'' probably refers to pyro-scapulimancy, see Maenchen-Helfen (1973), pp. 268-270, though the caution of Eisenberger (1938), p. 81, should still be noted. In modern Mongolia, the scapula was still scraped before it was burned (Montell [1945], p. 380).}\footnote[8]{Cooper (1936), p. 34; Eisenberger (1938), p. 50.}\footnote[9]{Satow (1879), p. 431; Eisenberger (1938), p. 81.}\footnote[10]{Rockhill (1900), pp. 187–188; cf. Wieger (1923), pp. 1676-1677.}\footnote[11]{Cooper (1936), p. 29.}\footnote[12]{On the hypothetical diffusion of pyroscapulimancy from the Old World to the New, see Speck (1935), p. 128; Cooper (1936), pp. 3031, 38; Hultkrantz (1968), p. 82; Needham (1971), p. 544. Pyro-scapulimancy was not, apparently, practiced by the Alaskan Indians on the Pacific watershed. If it came from Siberia, the culture trait was diffused in the area stretching from the Mackenzie and Peel rivers to the southeast, rather than south along the Alaskan coast. 1}
% source: scan 022, printed 5
twentieth centuries, it was reported from such areas as the south part of the Arabian
peninsula, Central Asia, Mongolia, Siberia, Manchuria, South and West China, Japan,
the Arctic, and Labrador.\footnote[13]{For the Arabian peninsula, see Andrée (1906), p. 155; Eisenberger (1938), p. 91. For Central Asia, Mongolia, and Manchuria, see Andrée (1906), pp. 147-150; \pinyinterm{ling-chunsheng} (1934), pp. 130-140; Eisenberger (1938), pp. 8189; Montell (1945), pp. 380-381; Bawden (1958); Roux and Boratav (1968), pp. 302-306. For Siberia, see Bogoras (1907), pp. 487-490; Andrée (1906), p. 151; Eisenberger (1938), pp. 78-79. For South and West China, see Lin Sheng (1963, 1964). For Japan, see Fujino (1960). For the Arctic, see Lot-Falck (1968), pp. 255-256. For Labrador, see Speck (1935), pp. 127-159; Hultkrantz (1968), pp. 74–83.} %13
 The animal which supplied the scapula (and sometimes other
bones) has varied; diviners have burned the shoulder blades of sheep, cattle, deer, reindeer,
saiga, caribou, seal, wild boar, pig, hare, dog, lynx, fox, beaver, camel, and porcupine.
But the attempt to seek reassurance about the future in burn cracks made in animal
shoulder bones satisfied a characteristic human need-in ways that historians of culture
have barely begun to explore—for well over five thousand years.
None felt that need more strongly, or elaborated scapulimancy more strikingly to
express their central concerns, than the kings of Bronze Age China. No matter who should
be given credit for the invention of the pyromantic art-sedentary Neolithic inhabitants
of North China, nomadic sheep herders of Central Asia, or members of a more primitive
circumpolar hunting culture\footnote[14]{Eisenberger (1938), pp. 101-106, summarizes various theories about the origins of scapulimancy and argues for a single point of origin. His own hypothesis (pp. 111-112) that apyromancy originated among hunters of the European Stone Age is unsubstantiated and, as he admits, fails to account for the presence of pyromancy east of the Caucasus. For more recent views, see Roux and Boratav (1968), p. 302, and Hultkrantz (1968), p. 83, n. 2. On the Chinese origins of pyro-scapulimancy, see \pinyinterm{zhang-bingquan} (1967), p. 851 and n. 119; Ho (1975), p. 285. It is not yet possible to specify the region (or regions) from which Neolithic pyro-scapulimancy may have first spread in China. An eastern origin has been suggested (\pinyinterm{li-ji} [1933], p. 576; \pinyinterm{shi-zhangru} [1954], p. 265), and scapulimantic remains have been excavated from 龙 and other Neolithic sites in Liaoning, \pinyintext{jilin}, Inner Mongolia, \pinyintext{hebei}, and \pinyintext{shanxi-province} (Itō [1962a], pp. 251, 259; \pinyinterm{zhang-bingquan} [1967], pp. 842844). At present, the earliest find is that from \pinyinterm{fuhegoumen} in Liaoning (see n. 3). The scapulimantic techniques of the \pinyinterm{qijia} culture in \pinyintext{gansu} are thought to have been more primitive than those found in the lower stratum at \pinyinterm{chengziyai} (\pinyinterm{liu-yuanlin} [1974], p. 103), but no conclusions about the origin and diffusion of the custom are possible until the relative age of all the sites and of the strata within them is known. Cf. n. 23.} %14
—the contribution of the \pinyinterm{shang} Chinese was threefold and
unique. First, they divined with turtle shells (plastromancy)\footnote[15]{According to Yang \pinyinterm{junshi-short} (1963), p. 192, n. 11, the term ``plastromancy,'' the pyromantic cracking of turtle plastrons, was first coined by \pinyinterm{li-ji}. In what follows, I frequently distinguish between shell (plastrons and carapaces) and bone (scapulas). Strictly speaking, the turtle shell was itself composed of bone (Carr [1952], p. 36); thus the general term ``oracle bone'' will be used to refer to both shells and bones when no distinction is intended. 11 1 ] 13 ] 0 ] 5 ] 0 ] 0 ] ] ] ]} %15
 as well as cattle scapulas,
thus introducing reptiles into what had been exclusively a mammalian realm. Second,
they cut orderly series of hollows or depressions into the bone or shell so that the omen
cracks were no longer free-form but structured, assuming specific shapes in specific places.
Third, they incised the divination records (known by the modern term \pinyinterm{jiaguwen},
``writings on shell and bone'') into the scapulas and plastrons, thus establishing themselves
% source: scan 023, printed 6
as the first fully literate civilization east of the Indus, the first consistent users of that logographic script which has characterized the writing of China and much of the Far East
ever since, the first people of China of whom true history may be written.
The account of \pinyinterm{shang} pyromantic procedures which follows primarily describes the
situation in the reign of king \pinyinterm{wu-ding} (period I);* the various changes in procedure, to
the extent that we understand their evolution with time, are treated in chapter 4. Much
of the discussion may, at first glance, appear forbiddingly technical. It is essential, however,
that we understand what the diviners did if we are to reconstruct and decipher the records
that they left behind, if we are to use those records with assurance as historical sources,
and if we are eventually to grasp what the diviners and the kings they served thought they
were doing. In attempting to understand a culture more than three thousand years distant
from our own, we need to use every available clue. The historiography of the oracle-bone
inscriptions involves more than just the inscriptions; it involves the scapulas and plastrons
as well, with their hollows and burn marks and cracks.
\section[The Divination Materials]{\origsecnum{1.2}The Divination Materials}
\label{sec:chapters-ch01:1.2}

Identification of the animals which supplied the bones and shells for \pinyinterm{shang}
pyromancy may tell us something of their habits and character, and thus of \pinyinterm{shang} patterns
of domestication. It may tell us something of their geographic distribution, and thus provide
information about the extent of \pinyinterm{shang} hegemony and of trade or tribute routes. It may
also assist modern attempts to duplicate \pinyinterm{shang} divination procedures and the subsequent
carving of inscriptions.
\subsection[Bovid and Other Bones]{\origsecnum{1.2.1}Bovid and Other Bones}
\label{sec:chapters-ch01:1.2.1}
Neolithic diviners used the scapulas of bovids (cattle and water buffalo), pigs,
sheep, and deer. \pinyinterm{shang} diviners continued to use these materials, but they relied increasingly
upon bovid bones, which, together with turtle plastrons (sec. 1.2.2), they were using almost
exclusively by the \pinyinterm{anyang} phase.\footnote[16]{For a useful list of the scapulimantic materials found in Neolithic and \pinyinterm{shang} sites, see Itō (1962a), pp. 251-252, and \pinyinterm{zhang-bingquan} (1967), pp. 842-848. Both articles overlook the early \pinyinterm{shang} scapula(s) (?) prepared with hollows, found at \pinyinterm{lusicun} in western \pinyintext{henan} in 1959 (KK [1964.9], p. 436, fig. 3.10). Unprepared scapulas have recently been found at \pinyinterm{erlitou}, \pinyintext{henan} (KK [1974.4], p. 241), \pinyinterm{cixian}, \pinyintext{hebei} (KKHP [1975.1], p. 98), \pinyinterm{ningcheng}, Liaoning (KKHP [1975.1], p. 126, plates 5.4-5), and \pinyinterm{qijia}, \pinyintext{gansu} (KKHP [1975.2], pp. 87-88). \pinyinterm{liu-yuanlin} (1974) discusses the various types of bone used and the increasing degree to which they were prepared by sawing, polishing, and hollow making (secs. 1.3; 1.5). See too n. 25.} %16
 The \pinyinterm{anyang} scapulas have been identified as those
of cattle (Bos exiguus Matsumoto) or water buffalo (Bubalus mephistopheles Hopwood), but
*For the five periods, \ref{sec:chapters-ch04:4.2} (see see sec. 4.2).
% source: scan 024, printed 7
the specific identifications should be treated with caution.\footnote[17]{The terms Bos spp. and Bubalus spp., indicating that specific identifications are still questionable, are probably to be preferred. The standard identifications given in the text, originally made by Matsumoto Hikoshichirō (1915), pp. 32-34, (1927), pp. 51-52, and Hopwood (1925), and revised by de Chardin and Young (1936), pp. 44-52, 56, 58 (cf. \pinyinterm{chen-mengjia} [1956], p. 5, quoting a 1953 letter from \pinyinterm{yang-zhongjian}, i.e., Young), are based not on the scapulas but on cranial fragments, mandibles, a large number of canon bones, and, above all, on the massive, sharply trihedral horn cores with their strongly inclined backward orientation. The link between these specimens and the plastromantic scapulas is circumstantial. Note that these identifications were made at a time when new scientific names were given to fossils with considerable freedom. Osteoarchaeologists have only started to deal with the many problems of postcranial identification. As with the case of the turtles (sec. 1.2.2), further study may permit the relegation of the specimens, especially the Bos, to synonymy with other known species. Note, for example, the comment of C. C. Young (1936), p. 515, n. 1, that a water buffalo ``of a type strikingly similar to the fossil Bubalus of China'' occurs as far west as Germany.} %17
 Furthermore, due to the
difficulty of distinguishing one type of scapula from the other, we cannot yet tell whether
the \pinyinterm{shang} diviners used scapulas from both types of animals or from just one.\footnote[18]{\pinyinterm{shi-zhangru} (1953), p. 9. The scapulas of cattle and buffalo may differ. According to MacGregor (1941), p. 444, a buffalo scapula in contrast to a cattle scapula has ``a convex anterior and a concave posterior border, as in the horse. The junction with the cartilage at the posterior edge shows a distinct notch. The acromium arises at an angle of 60° instead of 80 to 90°.'' And Richard H. Meadow tentatively concludes (letter of 7 October 1975), ``on the basis of one specimen of Bubalus here [the Museum of Comparative Zoology, Harvard University], it seems that it is relatively wider in relation to its length than is the case with Bos.'' But even if further study confirms these distinctions, it will be hard for the oracle-bone specialist to apply them since no \pinyinterm{shang} cartilage remains, the spine bearing the acromium was sawn off (sec. 1.3.1), and complete scapulas are rare.}
The preponderance of water buffalo remains over those of cattle, however, suggests that buffalo
scapulas would have been in greater supply.\footnote[19]{\pinyinterm{yang-zhongjian} and \pinyinterm{liu-dongsheng} (1949), p. 147, note that over a hundred cattle bones and over a thousand water buffalo bones had been found at \pinyinterm{anyang}.}%19
 Both animal types, which were probably
indigenous to North China and domesticated in \pinyinterm{shang} times, are now extinct.\footnote[20]{De Chardin and Young (1936), pp. 56, 58. \pinyinterm{zhou-mingzhen} and \pinyinterm{xu-yuxuan} (1957), pp. 460, 462, chart the distribution of fossil bovids in China; remains of Bubalus mephistopheles have been found only at \pinyinterm{anyang} and only during the Holocene. Cf. C. C. Young (1936), p. 515.}%20

The cranium, astragalus (ankle bone), and ribs of these bovids were also occasionally
divined at \pinyinterm{anyang}, but such bones were primarily used for practice inscriptions or record
keeping. Human and deer skulls were also inscribed with records, as were other bones
from deer and horses, but with the exception of one horse scapula, they do not appear to
have been used for pyromancy.\footnote[21]{\pinyinterm{rao-short} (1959), pp. 12-14, gives examples of each type; cf. \pinyinterm{gao-quxun} (1949), p. 155, n. 1, pp. 166, 184; \pinyinterm{chen-mengjia-spaced} (1956), pp. 5-6. Yang \pinyinterm{junshi-short} (1963), p. 195, n. 16, cites seven inscribed fragments of human skulls; \pinyinterm{li-yan-spaced} (1966), ), p. 2, refers to eleven. Pig and sheep scapulas and the leg bones of cattle or buffalo, all of which had been used in divination, have been found in small quantities at sites such as \pinyinterm{zhengzhou} and \pinyinterm{liulige}, but not at \pinyinterm{xiaotun} (焦, loc. cit.; cf. de Chardin and Young [1936], pp. 4452; Qu 万 [1961], p. 5; Itō [1962a], pp. 256-257). \pinyinterm{yan-yiping} (1955), p. 24, notes that \pinyinterm{luo-zhenyu} originally mistook the edges of fragmented scapulas to be rib bones. Some rib bones, however, appear to be correctly identified (e.g., \pinyinterm{chen-mengjia-spaced} [1956], plate 15, bottom; \pinyinterm{hu-houxuan} [1973], p. 60; \pinyinterm{zhou-hongxiang} [1976], pp. 1, 7-8, citing USB 10 and 11). For a period III horse scapula used in divination and inscribed, see KKHP (1958.3), pp. 70, 71. Gibson (1934), plate 8, reproduces a deer antler incised with divination inscriptions, but these are fake. \pinyinterm{dong-zuobin} (1961) refutes the earlier suggestion (cf. Britton [1935], pp. 1, 7) that the larger inscribed scapulas were those of elephants. 7 ] ] ] 0 ] ] 1 ] 1 1 0 ] 0 门 ] ] ] 0 ]} %21
 Most of the nondivination writings on bone are records
% source: scan 025, printed 8
of hunt or sacrifice from period V, often given with a \pinyinterm{shang} date, and may have been
carved on the bone of the animal or human victim.\footnote[22]{Typical examples are \pinyinterm{jiabian} 3939 (S244-2); 3940 (S537.2); 3941 (S244.2); \pinyinterm{jingjin} 5281 (S465.4, but adding yung , ``used in sacrifice,'' to Shima transcription); White 1915 (52.460), which \pinyinterm{xu-jinxiong} identifies as a carved and inlaid tiger bone, inscribed with a record of the tiger's capture (cf. Bone, pp. 96-98); \pinyinterm{yicun} 518 (S244.2), with ornamental carving on the other side.} %22

\subsection[Turtle Shells]{\origsecnum{1.2.2}Turtle Shells}
\label{sec:chapters-ch01:1.2.2}
The specific part of China in which a diviner first cracked a turtle plastron rather
than an animal scapula is not yet known,\footnote[23]{For what may be the earliest case of plastromancy, at \pinyinterm{gaohuangmiao} (\pinyintext{jiangsu}) in the Huai valley, see \pinyinterm{liu-yuanlin} (1974), p. 109; Itō (1962a), p. 252, however, classifies the bone and shell fragments as late \pinyinterm{shang}. For the original report, see KKHP (1958.4), pp. 12, 18, plate 8.1-4. in} %23
 but, in general, plastromantic remains have
been found only in association with \pinyinterm{shang} sites.\footnote[24]{According to \pinyinterm{shi-zhangru} (1954), pp. 267-268, pre-\pinyinterm{shang} plastromantic remains have been found only at Anshangcun in western \pinyintext{shandong} and Heigutui in eastern \pinyintext{henan}; Itō (1962a), pp. 251, 259, tentatively classifies these as late 龙 sites; \pinyinterm{zhang-bingquan} (1967), pp. 844, 846, classifies them as \pinyinterm{shang} sites. The issues are partly those of definition; and the sites may be transitional. It is also possible that the prepared plastron fragment found at Heigutui may have been intrusive from a later stratum (\pinyinterm{li-jingtan} [1947], p. 118).} %24
 Crudely prepared turtle shells have been
found in the lowest \pinyinterm{shang} strata at \pinyinterm{zhengzhou}; they become increasingly sophisticated
in their preparation as they become more numerous by comparison with scapulas.\footnote[25]{For Neolithic and \pinyinterm{shang} plastromantic materials, see the articles cited in n. 16. New reports of excavations carried out just south of \pinyinterm{zhengzhou} in 1955 indicate that unprepared plastrons were found in the \pinyinterm{nanguanwai} stage; plastrons with chiseled-out hollows were found in strata corresponding to the Lower and Upper \pinyinterm{erligang} stages. But in every case bovid scapulas predominated (KKHP [1973.1], pp. 70, 79, 88, 91, plates 3.1, 4.2, 13.8). It is only in the \pinyinterm{renmin-gongyuan} phase that plastrons become more numerous than scapulas (\pinyinterm{zhengzhou} \pinyinterm{erligang}, p. 38). There have also been recent reports of plastromantic remains associated with the \pinyinterm{shang} cultural sites at \pinyinterm{taixicun} in \pinyintext{hebei} (KK [1973.5], pp. 267268) and at Qiuwan in \pinyintext{jiangsu} (KK [1973.2], p. 74). In both cases, the plastrons had been prepared for divination; prepared bovid scapulas were also found, and, in \pinyinterm{taixicun} at least, were more numerous than the plastrons.} %25
 It
has been suggested that \pinyinterm{shang} pyromancy represents the coalescence of two Neolithic
traditions: the first, a scapulimantic tradition, found along both sides of the Wei and Yellow
rivers and to their north; the second, a plastromantic one, centered south and east of the
Yellow River.\footnote[26]{\pinyinterm{shi-zhangru} (1954), pp. 268–269, proposes that scapulimancy, found in the region from \pinyintext{shandong} and \pinyintext{henan} to \pinyintext{liaodong} (a region which tradition assigns to the Dongyi) may be associated with the black-pottery culture (cf. \pinyinterm{li-ji} [1934], p. xv); and that plastromancy, originating in the Huai valley (the region of the Huaiyi) may be associated with the impressed-pottery culture. See too, Itō (1962a), p. 259, map; cf. Kaizuka (1967), p. 197; Ho (1975), p. 283, n. 33; also n. 23 above. Chinese myth, incidentally, credits the culture hero Fu Xi with the invention of plastromancy, but the attribution first appears in Wuyuan of Ming times (\pinyinterm{ling-chunsheng} [1971], p. 22). Nor may any credence be given to Sung accounts of a divine tortoise shell sent from the region of present Vietnam to the court of the sage emperor Yao in 2361 B.C. (Thai Van-kiem [1963], pp. 584-585), though a southern origin for both the custom and some of the \pinyinterm{shang} shells cannot be excluded (see n. 43).} %26
 Further research and excavation may show whether this hypothesis, not
yet well supported by Neolithic finds, is true.
% source: scan 026, printed 9
As for the types of turtle used, a thorough paleontological examination of the inscribed
shells, long needed, has now been undertaken by Berry, whose methods and conclusions
are given in appendix 1. Working with the reconstructed plastrons in \pinyinterm{jiabian}, \pinyinterm{pinbian},
and \pinyinterm{yibian}, he finds four chelonian species represented: 27 \listitem{1} Ocadia sinensis, \listitem{2} Chinemys
(=Geoclemmys) reevesi, \listitem{3} Mauremys (=Clemmys) mutica, and \listitem{4} a single example, Testudo
emys. The \pinyinterm{anyang} field tortoise, Testudo anyangensis, who still crawls through many monographs, is probably not extinct, but taxonomically superfluous; the \pinyinterm{shang} diviners may
be absolved of the charge of exterminating the species. 28
A list of Berry's identifications is provided in tables 3 and 4. The following conclusions
may be drawn from his research. Ocadia sinensis, Chinemys reevesi, and Mauremys mutica are
all aquatic, freshwater turtles; they are generally thought good to eat; and they can be
kept in captivity. The fact that Ocadia sinensis and Mauremys mutica are now found only in
South China raises the possibility that their shells may have been among those which we
know, from the marginal notations, to have been imported from the south (secs. 1.2.4
and 1.4). Further research may succeed in linking specific turtle varieties with the names,
carved on the shell, of specific statelets or regions in the \pinyinterm{shang} sphere of influence.29 The
unique tortoise Testudo emys (\pinyinterm{pinbian} 184) is of a species now found in Southeast Asia and
the Malay Archipelago; if its shells traveled to North China in \pinyinterm{shang} times, they did so
in small numbers.
It may also be remarked that the turtles used in divination were probably female.
This is still the custom in Taiwan today and may be explained by the fact that the plastron
of the female is smoother than the plastron of the male. Its more uniform thickness makes
it more suitable for pyromantic cracking, its greater smoothness more suitable for the
recording of inscriptions.\footnote[30]{On the basis of plastromantic custom in present-day Taiwan, \pinyinterm{zhang-guangzhi} (1972), p. 22, has suggested that the \pinyinterm{shang} used female turtles. (胡 徐 [1782-1787], \ref{ch:4} (see ch. 4), p. 3b, reporting on \pinyintext{jiangsu} plastromancy in the first part of the eighteenth century, refers to distinguishing the turtles by sex, but he does not indicate that one was preferred.) The presumed preference for females is supported by Berry (letter of 6 October 1975): ``It is very likely that all (or nearly all) of the plastra used for divining were those of female turtles. The plastron of male Ocadia and Chinemys is rather deeply concave and would presumably be more difficult to inscribe. There would also probably be a difference in the way the 'inksqueezes' would look if males were used.'' Female plastrons also tend to be larger than male ones (Mao [1971], pp. 34–52). ] ] ] 0 0 ] 1 ] 门 ] 1 ] ]} %30

27. It is possible that different turtle types may
be represented by different oracle-bone graphs.
Ting Su (1969) and \pinyinterm{ling-chunsheng} (1971),
pp. 27-29, discuss the various \pinyinterm{shang} graphs; see
too, \pinyinterm{xu-zhongshu} (1931), pp. 530-531.
28. E.g., Crump (1963), p. 170; Vandermeersch (1974), p. 38.
29. For example, all the marginal notations
``Que sent in 250 (shells)'' (appendix 3,
n. 11) were inscribed on Chinemys reevesi plastrons.
The location of the statelet of Que is in dispute.
\pinyinterm{hu-houxuan} (1944), p. 55, and Shima (1958),
p. 464, locate it to the west of the \pinyinterm{shang}; 陳
\pinyinterm{mengjia-short} (1956), p. 298, locates it to the southwest; Itō (1962a), p. 268, to the southeast; Ting
Shan (1956), p. 125, locates it to the south of
\pinyinterm{anyang} in the \pinyinterm{kaifeng} region; Shirakawa
(1957), pp. 14, 33, also locates it south of the
capital, two to three days' journey on foot. For a
detailed study of the geography, see ibid., pp. 54-
57. Unfortunately, none of the inscriptions indicating a southern source (see n. 43 below) are on
plastron fragments large enough to be identified.
% source: scan 027, printed 10

\subsection[The Ratio of Bone to Shell]{\origsecnum{1.2.3}The Ratio of Bone to Shell}
\label{sec:chapters-ch01:1.2.3}

The editors of more recent collections of oracle-bone rubbings generally indicate
whether a particular fragment is of scapula bone or of turtle shell.\footnote[31]{Liu-lu (1945) was the first collection in which shell and bone were so identified. The fragments in \pinyinterm{renwen}, Menzies, and White are arranged by period and then by divination material, with an identifying ``B'' (bone) or ``S'' (shell) appended to the rubbing number; the ``B'' and ``S'' notations are used in \pinyinterm{minghou}; \pinyinterm{xu-jinxiong}'s commentaries to Menzies and White identify the particular area of scapula or plastron from which a fragment came. The registration numbers given for the fragments in \pinyinterm{jiabian}, \pinyinterm{zhuihe}, \pinyinterm{pinbian}, and \pinyinterm{yibian} (8532-8637 lack registration numbers) also identify the material as shell or bone. Here a word of explanation may be in order. The registration number, which was recorded after excavation on the bone or shell fragment itself, is to be distinguished from the rubbing number printed under, and generally used to refer to, the published inscription (see \ref{ch:3} (see ch. 3), n. 5). Thus, the fragment published as \pinyinterm{yibian} 7574 (the rubbing number, printed in larger, heavier type) was first assigned the registration number 13.0.16015; \pinyinterm{jiabian} 2906 had the registration number 3.2.0877. Both registration (or, in some cases, accession) and rubbing numbers are given beneath the rubbings in \pinyinterm{jiabian}, \pinyinterm{zhuihe}, \pinyinterm{pinbian}, \pinyinterm{yibian}, Menzies, and White. The registration number used by the Academia Sinica is tripartite: the first number records the dig number; the second identifies the type of object (o = inscribed shell; 1 = uninscribed shell; 2 = inscribed bone; 3 = uninscribed bone; on the basis of \pinyinterm{jiabian} 3939, 6 = inscribed bovid skull); the third is the identification number assigned to the object when it was excavated (董 [1949], p. 2). Thus, \pinyinterm{yibian} 7574 was excavated in the thirteenth dig, it is a piece of inscribed shell, its identification number was 16015; \pinyinterm{jiabian} 2906 was excavated in the third dig, it is a piece of inscribed bone, its identification number was 0877.} %31
 For the older collections,
identification may be made from the rubbing on the basis of fragment shape, suture lines,
and so forth, though in some cases there are no visible clues.
For the historical period as a whole, turtle plastrons and bovid scapulas were probably
used by the \pinyinterm{shang} diviners in about equal numbers, though their relative popularity may
have varied with time. (For differing estimates and the technical issues involved, see
appendix 2.) It is possible that there was a preference for bone in period III + IV, for
example, and for shell in period V.\footnote[32]{According to \pinyinterm{chen-mengjia-spaced} (1956), p. 142, the early and late periods (i.e., periods I-IIIa and V) used both bone and shell, periods IIIb and IV generally used bone and rarely used shell. \pinyinterm{xu-jinxiong} (1973a), pp. 2-3, 48-49, finds that period III inscriptions which record the diviner's name were mostly on shell, and that all period IV inscriptions (with a few exceptions in IVb) were on bone. He also finds that shell was most often used in periods I, II, and V ([1973], pp. 26, 126). On the preference for shell in period V, see Itō (1962a), p. 253. Table 5 does not fully support these views; and the apparent preference for scapulas in periods III + IV or IV may reflect the accidents of discovery rather than a real preference (see sec. 4.3.3.3). On the use of carapaces for period V sacrifice divinations, see n. 57.}%32
 The fact that the diviners of the Royal Family group
(sec. 2.2.1.1) showed a preference for shell suggests that they may have valued plastrons
and carapaces more than scapulas as a divination medium.\footnote[33]{The distinction may be seen from pit E16 which produced 243 rubbings on shell and 20 on bone; the Royal Family group inscriptions were all on shell. The contents of pit F36, a pure RFG find, consisted entirely of plastron fragments. Virtually all the RFG inscriptions from pits B119, YH6, YH 127 were inscribed on shell; the pure RFG find, YH344, produced twice as many shell rubbings as bone rubbings. Of the 261 RFG rubbings given in \pinyinterm{renwen}, 241 are on shell, only 20 on bone. Adjustment according to the fragmentto-whole-bone ratios given in appendix 2, sec. 2, would not invalidate the significance of these figures. On the other hand, a majority of the Menzies inscriptions which \pinyinterm{xu-jinxiong} dates to period IV are on bone; a significant proportion of these should perhaps be dated to period IIa (see \ref{ch:4} (see ch. 4), n. 44), and some to the RFG (\ref{ch:2} (see ch. 2), n. 18). But the matter needs more research and will be discussed in Studies. 1}%33
 But, generally speaking, no
For further discussion of these registration numbers, see \ref{ch:4} (see ch. 4), n. 219.
% source: scan 028, printed 11
evident principle governed the use of bone or shell, the choice of which appears to have
been unrelated to the matter being divined.\footnote[34]{Later research has not supported \pinyinposs{dong-short} early suggestion ([1929b], pp. 208-209) that shell was used exclusively in the first stage of pyromancy. He subsequently proposed ([1945], pt. 2, ch. 10, p. 1b) that divinations about the night used plastrons in periods III and V, carapaces in periods I, IV, and V, and scapulas in periods I and II. But these generalizations appear to be qualified by \pinyinterm{jiatu} 202, a period III carapace used for night divinations (for the meaning of the italicized rubbing number, see \ref{ch:4} (see ch. 4), n. 1), and by the situation in \pinyinterm{renwen} (1429-1458, 1623) and Menzies (1512-1525), where the great majority of the period II night divinations used shell. 董 also suggested ([1965], p. 117) that the New School diviners (on the Old and New Schools, \ref{sec:chapters-ch04:4.2} (see see sec. 4.2)) had rules about using bone or shell for particular topics. He does not give examples, and I have not found supporting evidence. The rule would appear to be contradicted by the case cited in n. 35. For further discussion of bone, shell, and divination topic, see appendix 2, sec. 3.}%34
 The diviner's indifference to the medium
with which he worked, at least for certain periods, is well demonstrated by a period V
case in which a series of divinations about the same event was recorded on a series of scapulas
and plastrons.\footnote[35]{In period Vb (\pinyinposs{dong-short} New School), the campaign against the Renfang was recorded on both bone and shell; \pinyinterm{dong-short} (1945), pt. 2, \ref{ch:5} (see ch. 5), pp. 18b-21a, reconstructs three of the plastrons and seven of the scapulas involved; cf. ibid., ch. 8, p. 8b; ch. 9, pp. 59a-60a.}%35
 Scapulas and plastrons varied greatly in size, but no relationship has been
established between the size of the bone or shell and the importance of the topic.\footnote[36]{The largest plastron found (\pinyinterm{pinbian} 184) is that of the tortoise, Testudo emys; it is 43.5 cm. long. The shortest I have noted (\pinyinterm{yibian} 2903) has been identified as Chinemys reevesi; this turtle plastron is 11.7 cm. long. For traditional accounts of plastron size, see \pinyinterm{ling-chunsheng} (1971), pp. 24, 33, 34. The largest scapula I have noted is \pinyinterm{ninghu} 1.110, approximately 43 cm. long; the smallest is Menzies, plate 275, which may originally have been approximately 23 cm. long. On the matter of size, cf. Kaizuka (1967), p. 199; Shirakawa (1972), p. 19.} %36


\subsection[The Sources of Bone and Shell]{\origsecnum{1.2.4}The Sources of Bone and Shell}
\label{sec:chapters-ch01:1.2.4}

It thought that the scapulas used by the \pinyinterm{shang} diviners came in part from the
cattle or water buffalo whose slaughter supplied the \pinyinterm{shang} sacrifices (and the royal table)
and that many of the animals came from the \pinyinterm{shang}'s own herds.\footnote[37]{\pinyinterm{hu-houxuan} (1944), pp. 55a, 59b-60a; (1944a), p. 6a; Itō (1962a), p. 255. On the gregarious and generally tractable nature of water buffalo, see, for example, Walker et al. (1968), p. 1423; on their edibility, see Levine (1919), p. 37. The fact that no water buffalo skeletons, but only those of cattle, have been found in the burial pits associated by modern archaeologists with sacrifice (\pinyinterm{shi-zhangru} [1953], p. 9) and the further fact that, to the best of my knowledge, none of these cattle skeletons had had their scapulas removed suggest that if the pyromantic scapulas were taken from sacrificial victims, they would have come from victims that were dismembered for the sacrifice, with part of the flesh offered to the spirits, part consumed by the worshippers, and the scapulas used to divine the spirits' reactions. Since several sacrifice words appear to involve dismemberment (e.g., mao p [S507.1-508.3]; see Serruys [1974], p. 92, n. 6), this seems a plausible speculation.} %37
 Bovids, presumably on
the hoof, were occasionally sent by other statelets,\footnote[38]{E.g., \pinyinterm{pinbian} 81.5–8.} %38
 and marginal notations on scapula
sockets or on bone surface (sec. 1.4) record the non-\pinyinterm{shang} provenance of certain scapulas.\footnote[39]{\pinyinterm{hu-houxuan} (1944a), pp. 35a-44b. It is generally assumed that scapulas and plastrons were sent as tribute; \pinyinterm{zhou-hongxiang} (1970/71), pp. 362-364, estimates that there were about 115 tribute-paying areas. There is also the possibility that some of the shipments involved trade or expropriation, however, and thus the marginal notations might have been accounting rather than, or as well as, tribute records. On the difficulties involved in characterizing natural resources flow on the basis of archaeological evidence, see \pinyinterm{zhang-guangzhi} (1975). II ] 0 ] ] ] ] 1 ] n ] 1 0 ]}³9
% source: scan 029, printed 12
The report of the \pinyinterm{xiaotun} villagers that they once found a store of several hundred
turtle shells of all sizes, large and small, with plastrons and carapaces complete,\footnote[40]{\pinyinterm{dong-short} (1929b), p. 212.} %40
 suggests
that turtles were reared, or at least tended, in captivity. Chinemys reevesi, in particular, was
probably native to the region, and all the turtles identified by Berry can be kept in captivity
(see appendix 1, sec. 4). Since the diviners probably used more than one plastron a day
(appendix 3, sec. 4), the \pinyinterm{shang} may have engaged in turtle farming to supply divination
materials as well as food.
A significant proportion of the turtle shells used for plastromancy, however, appears
to have come from outlying regions. The marginal notations (sec. 1.4) reveal that turtles
were sent in, in the same way as were the horses, dogs, bovids, and 羌 captives recorded
in the divination charges.\footnote[41]{\pinyinterm{zhang-bingquan} (1967), pp. 852-853. See, for example, \pinyinterm{pinbian} 157.11-12; 342.4-5; 342.2-3; 81.5-8;574-} %41
 \pinyinterm{shang} notations record the receipt of large numbers of turtle
shells with some cases involving 250, 300, 500, or even 1,000 shells, at one time.\footnote[42]{\pinyinterm{hu-houxuan} (1944), p. 55b. At least six \pinyinterm{yibian} plastrons (S352.2) record the ``taking'' (on this word, see n. 62) of 1,000 shells; the fact that the five that can be identified were all Chinemys reevesi and were all excavated from the same pit, YH127, suggests that they were part of a single shipment. See appendix 3, nn. 12, 14.} %42
 The exact
location of the donor statelets may never be fully determined, but notations occasionally
record that shells came from the east, west, and south.\footnote[43]{\pinyinterm{hu-houxuan} (1944a), pp. 6b-8a. \pinyinterm{qianbian} 4.54.6 and \pinyinterm{kufang} 592 (D) (both \$246.2) may refer to turtles from the west. For the names of the tribes which sent in turtles and which are thought to have been situated to the east or west of \pinyinterm{shang}, see 胡 (1944), p. 55b, but cf. n. 29 above. On the southern origin of the shells, see too \pinyinterm{shi-zhangru} (1954), p. 269. \pinyinterm{pinbian} 621.6 and \pinyinterm{qianbian} 4.54.4 and 4.54.5 (both S246.2) refer to turtles from the south. For clues about the political and cultural situation south of \pinyinterm{shang} and the possible existence of friendly tributaries in that region, see Kane (1974/75), p. 85 and n. 13; Hayashi Minao (1971), pp. 43-45 and n. 67, whom Kane cites, provides a map of the archaeological situation. Further insights may be gained from the excellent maps in Barnard and Satō (1975), especially map 1a, p. 92.} %43
 A study of the provenance of turtles,
as recorded in pre-\pinyinterm{han-dynasty} and later texts and in local gazetteers, confirms that in traditional
times turtles were primarily found in the Yangtze region, the \pinyintext{sichuan} area, and the
south.\footnote[44]{\pinyinterm{hu-houxuan} (1944a), pp. 8a-9a. Cf. \pinyinterm{ling-chunsheng} (1971), pp. 22-24.} %44
 As we have already seen (sec. 1.2.2 and appendix I, sec. 4), Ocadia sinensis and
Mauremys mutica now have a southern habitat. Thus \pinyinterm{shang} records, traditional accounts,
and zoological studies support the view that \pinyinterm{shang} plastromancy depended in part upon
imported shells for its raw material.

\section[Initial Preparation]{\origsecnum{1.3}Initial Preparation}
\label{sec:chapters-ch01:1.3}

The bone or shell surface was usually worked in a variety of ways before it was
scorched by the diviner's brand.\footnote[45]{A few inscribed scapula fragments whose surface had not been scraped smooth were excavated from pit F37 (\pinyinterm{jiabian} 430; 435-437; and 1.2.0197; see 董 [1929b], p. 209, and the relevant \pinyinterm{jiabian} k'ao-shih entries), but they are rare in the historical period (see secs. 4.3.2.2 and 4.3.2.3). 鮮} %45

% source: scan 030, printed 13

\subsection[Scapulas]{\origsecnum{1.3.1}Scapulas}
\label{sec:chapters-ch01:1.3.1}

\subsection*{Initial Preparation}

In the case of bovid shoulder blades (fig. 1), the scapular cartilage was sawn away,
leaving only the bony part of the scapula; the spine on the back (or lateral) surface of the
bone was cut away up to the socket so that the rear surface was leveled and roughly parallel
with the front surface.\footnote[46]{"Back'' is the term used by modern oraclebone specialists to refer to the surface of bone (or shell) into which the hollows (sec. 1.5) were generally bored (for exceptions, see n. 77). In the case of scapulas, this was usually the lateral side where the bony ridge had been removed (``lateral” is the anatomical term, indicating that side of the scapula which, in the living animal, faced outward). The cracks and inscriptions generally appeared on the ``front'' (or ``costal'') surface of the scapula. In life, the socket was situated at the bottom of the bone, but since the \pinyinterm{shang} inscriptions must usually be read with the socket at the top (for exceptions, see \ref{ch:2} (see ch. 2), n. 115), oracle-bone scholars refer to the socket-end as the top, and the broad, flat edge as the bottom. Scholars also refer to the side of the scapula on which the socket had been notched as the ``inner'' edge, as opposed to the ``outer'' edge on the unnotched side (\pinyinterm{xu-jinxiong} [1974], pp. 209, 254-255). For the anatomy of the scapula (fig. 1), see Sisson (1953), pp. 84-86, 126, 145–146.} %46
 From one-third to one-half of the socket (the glenoid cavity) was
sawn away on the rear side, leaving a half-moon-shaped section on which the marginal
notations (sec. 1.4) were frequently inscribed. The tuberosity\footnote[47]{This consisted of the tuber scapulae and the coracoid process to which tendons and muscles were attached; probably the gristle was hard to clean, and it was found easier to cut the bone away entirely.}%47
 near the edge of the socket
was also sawn away so that a right-angled corner appears to have been cut out of the
bone.\footnote[48]{\pinyinterm{dong-short} (1954), p. 458, fig. 1. Presumably, the graph (S303.1), read as chan, ``to read the cracks, prognosticate'' (\ref{sec:chapters-ch02:2.7} (see see sec. 2.7)), depicted the scapula notched in this way. When the front side of an actual scapula is facing the viewer and the missing angle is in the top right position, then modern scholars refer to it as a right-hand scapula; if it is in the top left position, then it is a left-hand scapula (\pinyinterm{chen-mengjia-spaced} [1956], p. 10). In the case of unprepared bones, when the socket is at the top and the spine (facing the viewer) is on the left side, then the scapula is termed a right scapula, and vice versa (cf. Bawden [1958], p. 7).} %48
 The neck of the socket was also leveled (fig.\footnote[49]{The spine and socket were occasionally left intact during the historical period; such unprepared scapulas may have been used for calligraphy practice (\pinyinterm{zhou-hongxiang} [1976], p. 8, citing the ``unusual, if not unique'' case of US 40). \pinyinterm{liu-yuanlin} (1974), p. 128, concludes that prior to period I most scapulas were unprepared in this way.} 2).49
It should be noted that bones and shells were occasionally cut into oblong tablets.\footnote[50]{E.g., \pinyinterm{jiabian} 2244–2247; 2870.}%50

The number of such regular-sided tablets is larger than would appear from rubbings
because the corners have been broken off, or the rubbing corners, lacking characters, have
not been reproduced.\footnote[51]{\pinyinterm{liu-yuanlin} (1969). The function of such tablets, which frequently record period II divinations associated with the king, has not yet been determined.}%51


\subsection[Turtle Shells]{\origsecnum{1.3.2}Turtle Shells}
\label{sec:chapters-ch01:1.3.2}

Once the turtle had been killed and the shell cleaned of meat and cartilage, the
dorsal portion (the carapace) and the ventral portion (the plastron) were sawn apart at the
bridge, which remained attached to the plastron.\footnote[52]{The suggestion of Ting Su (1969), p. 107 (cf. \pinyinterm{ling-chunsheng} [1971], p. 28), that \pinyinterm{fujia-title} 2 (S464.3) refers to the cleaning of the shell is 13 [ ] ]] ] 11 1}%52
 The rising edge of the bridge was then
% source: scan 031, printed 14
trimmed and the outer circumference of the entire plastron was smoothed.\footnote[53]{\pinyinterm{dong-short} (1929), p. 84, fig. 8. For the technical vocabulary, both Chinese and English, used to describe the various segments of the shell, see 并 知 (1931), \pinyinterm{dong-short} (1958), pp. 909-911; cf. Malcolm A. Smith (1931), facing p. 50; Carr (1952), pp. 37-39; Romer (1956), p. 32. The terms used in the text are shown in fig. 3.}\footnote[54]{For the number of shields and plates, see Malcolm A. Smith (1931), p. 50. Fig. 3 shows the differing pattern of outer scutes and inner bony plates (cf. appendix 1, sec. 2). The rubbings listed 'in tables 3 and 4 all reveal the scute pattern to a significant degree. For plates with the scutes stripped away, see, e.g., \pinyinterm{pinbian} 280; 289; \pinyinterm{yibian} 3426.}\footnote[55]{\pinyinterm{dong-short} (1929), pp. 83–87; \pinyinterm{chen-mengjia} (1956), p. 10; \pinyinterm{zhou-hongxiang} (1969), p. 35. The hollows (sec. 1.5) were generally cut into the back or interior side of the shell. The inscriptions were written on the front, exterior side. Modern scholars refer to the front, head-end of the plastron as the top and to the rear, tail-end as the bottom. Right and left are defined from the viewpoint of one facing the front, exterior side of the plastron (cf. \pinyinterm{zhang-bingquan} [1956], pp. 229–230).}\footnote[56]{\pinyinterm{chen-mengjia} (1956), figs. 2–5; \pinyinterm{zhou-hongxiang} (1969), figs. 16-18. Carapaces that were not sawn in this way are rare (\pinyinterm{yan-yiping} [1967], pp. 6-8, reconstructs one example). \pinyinterm{dong-short} (1949a), p. 265 (cf. [1945], pt. 1, \ref{ch:3} (see ch. 3), pp. 3a-b), suggests that the carapaces were split in two to form pairs like right and left scapulas.}\footnote[57]{\pinyinterm{dong-short} (1945), pt. 1, \ref{ch:3} (see ch. 3), pp. 3a-b, 8a, 12a, 18b; pt. 2 \ref{ch:2} (see ch. 2), p. 1b. Some of the RFG divinations also appear on carapace (\pinyinterm{li-xueqin} [1958], pp. 61, 68; cf. 董 [1945], pt. 2, \ref{ch:4} (see ch. 4), p. 4b); see, in particular, the large number of carapace fragments (reproduced in \pinyinterm{yibian} 1-237) found in pit B119.}\footnote[58]{\pinyinterm{zhang-bingquan} (1954), p. 225.}
% source: scan 032, printed 15

\subsection[Polishing]{\origsecnum{1.3.3}Polishing}
\label{sec:chapters-ch01:1.3.3}

some of the bones and shells found at \pinyinterm{xiaotun} were as highly polished as jade.\footnote[59]{\pinyinterm{dong-short} (1929b), p. 193; \pinyinterm{zhang-bingquan} (1954), p. 225. Cf. Creel (1935), p. 176, who notes that the bones were ``often scraped and polished so that they still gleam like a metal mirror.'' It is thought that scapulas were polished with particular care in periods I and IIa (\pinyinterm{liu-yuanlin} [1974], p. 118).}%59
 The
care lavished on their preparation suggests the high value the \pinyinterm{shang} placed on them.

\section[Marginal Notations]{\origsecnum{1.4}Marginal Notations}
\label{sec:chapters-ch01:1.4}

After the bones or shells had been sawed, scraped, and smoothed but before the
hollows (sec. 1.5) were bored or chiseled, marginal notations, which recorded such matters
as the source and number of shells and bones sent in, the persons associated with their
handling, and the day-date, were frequently made.\footnote[60]{This describes the situation under \pinyinterm{wu-ding}; marginal notations were little used after period I (see sec. 4.3.1.6). The basic studies are \pinyinterm{hu-houxuan} (1944, 1944a, and 1948); see too \pinyinterm{dong-short} (1954). The fact that marginal notations were partly erased by the hollows (as in \pinyinterm{pinbian} 586; \pinyinterm{yibian} 5270; 7491) indicates that the notation must have been carved before the hollows were made (董 [1954], p. 467; \pinyinterm{chen-mengjia} [1956], p. 18).} %60
 On turtle plastrons, these administrative records were carved on the back of the bridge (and are thus called by the modern
term, chia-ch'iao k’e-tz'u ), on the xiphiplastron (see fig. 3), or on the lower part
of the carapace, near the spine. On scapulas, they might be carved either in or near the
socket itself (hence ku-chiu k'e-tz'u ) or on the lower part of the bone, either front
or back, near the edge.\footnote[61]{\pinyinterm{dong-short} (1949a), pp. 267–268. \pinyinterm{chen-mengjia-spaced} (1956), p. 13, points out that these areas where the notations were carved on the plastrons were well removed from the pyromantic cracking. He also notes that the marginal notation on Ninghu 1.111 (D) (which he dates to period IVa) is inscribed upside down.}%61

The ``income notations which used the verbs ju \textasciicircum{} (``[cause to enter=] send in''),
lai (``[cause to come= =] send''), and '' (“bring, lead\footnote[62]{\pinyinterm{hu-houxuan} (1944a), pp. 1ob ff., attempts to distinguish between the meanings of ju (S256.1-257.1) and lai (S200.2); for (e.g., \$7.4-8.1), see Takashima (1973), p. 370, n. 28; Serruys (1974), p. 98, n. 13, summarizes the various possibilities: ``bring along, lead (armies, men), bring (sacrificial victims, etc.), present (as offering, gift), to hold, use (something)." ?]'') 62 indicate that large numbers
of shells 250, 300, 500, and so forth- -were received at one time.\footnote[63]{\pinyinterm{hu-houxuan} (1944), p. 55b. It has been suggested that the larger payments may have referred to tribute shipments from the south (\pinyinterm{renwen} shakubun, pp. 353-354). 胡 (ibid., n. 77 and [1944], pp. 11b-12b) has found notations which in total recorded the coming of 12,334 shells. This figure is deceptive, however, since two notations recording the arrival of 300 shells are more likely to be duplicate notations recorded on two shells belonging to the same 300-shell shipment rather than notations recording two separate shipments totaling 600 shells; on this point, see appendix 3, n. 12. 15 1 B 1 ] D 1 11 1] ]} %63
 In addition to indicating
the size of the shipments, the figures are significant because they help us estimate the extent
to which our present store of divination inscriptions is incomplete (appendix 3, sec. 2).
Since the ju and lai notations generally involve large numbers of turtle shells, and since
they never include a kan-chih day-date, it has been suggested that these entries refer to shipments sent from afar. It may be that the verbs ju and lai never appear in the scapula notations because cattle or water buffalo were not driven to the \pinyinterm{shang} capital from great
% source: scan 033, printed 16
碎五三父下的三甲\textasciitilde{}\textasciitilde{}

\subsection*{Marginal Notations}

distances. By contrast, notations using the phrase ch'i tzu = \# (S115.4-116.1), (``request
from, requisition from\footnote[64]{There are many theories about the meaning or meanings of the graph =. Here I follow the majority of scholars in taking it as antecedent to ch'i . In this context it appears to have the sense of ``request'' or ``take.'' For a good resumé of the scholarship, see \pinyinterm{renwen} 949, shakubun, p. 329, n. 1; CKWT, pp. 151–158. ?]''), 64 record far smaller numbers of bone or shell and usually
contain the day-date. This suggests that such records refer to consignments requested or
requisitioned from nearby localities, allowing the date the bone or shell was obtained to
be established with precision, as in the socket notation “On \pinyinterm{guihai}, Pei (?) requisitioned
from Yü ten (scapula) pairs. (Recorded by)\footnote[65]{This notation has been found on two bones: \pinyinterm{cuibian} 1503 (S170.4) and \pinyinterm{nanbei}, ``Ming'' 7 (D) (S115.4) (=\pinyinterm{minghou} 1761c), though the second lacks the name of the recorder; presumably, there were originally eight more such notations. The graph (S151.3-152.3; 482.4-483.2) appears only in the marginal notations of scapulas or carapaces. It is thought to refer either to a pair of right and left scapulas or to the pair of right and left sections into which the carapace was divided (e.g., \pinyinterm{pinbian} 65; \pinyinterm{yibian} 3330; 5281; cf. n. 56). The plastron, since it was used whole, was not recorded as a 1 (\pinyinterm{hu-houxuan} [1944], pp. 64a, 68a; \pinyinterm{renwen} 1091, shakubun, p. 351, n. 3; Qu 万 [1965], p. 75). The graphs ((the half-moon shape of a scapula socket?) or (a picture of a scapula) apparently referred to a single scapula, i.e., to half a scapula pair (\pinyinterm{guo-moruo-spaced} [1934], pp. 9-10; cf. Huang Caijun [1964], p. 433), as in ``Fu Xi ritually prepared seven (scapula) pairs and one half-pair. (Recorded by) Pin'' (\pinyinterm{houbian} 2.33.10 [S151.4]).} ``65 Notations on bone using the verb =
outnumber those on shell by four to one, suggesting further that the \pinyinterm{shang} tended to take
scapulas rather than turtle shells from regions close to their capital.\footnote[66]{\pinyinterm{renwen} shakubun, pp. 353, 354; Itō (1962a), pp. 254-256.} %66
 Distinctions such as
these may eventually throw light on the political economy of the \pinyinterm{shang} state.
The other verb frequently found in marginal notations is T (S151.3-152.3). Its
meaning in this context has been much debated.\footnote[67]{Among the meanings proposed for T in this context are: a ritual or sacrifice with regard to the bone or shell; the making of the hollows (see n. 71); the sending in or contributing of bone or shell; the placing or taking of the bone or shell. For a summary of the scholarship, see \pinyinterm{hu-houxuan} (1944), pp. 60a-b; \pinyinterm{renwen} shakubun, p. 351; Itō (1962a), pp. 253-255, 267, n. 92. See too n. 74 below.} %67
 Since the activity it represented was
performed prior to the pyromantic burning,\footnote[68]{This is made clear by \pinyinterm{yibian} 3330 (S412.4) and 5281, a carapace pair, both with the notation ``钱 (?) sent in 7; Fu 京 ritually prepared. (Recorded by) 韦.'' None of the hollows on either piece was burned (\pinyinterm{chen-mengjia-spaced} [1956], p. 18). For classical passages about rituals to prepare the turtle shell for divination, see \pinyinterm{hu-houxuan} (1944), p. 63; Kaizuka (1967), pp. 199–200; Shirakawa (1972), p. 19. The original passages may be found in SKK, ch. 128, pp. 26-27; CLCY, ch. 48, pp. 2a-b (Biot [1969], vol. 2, p. 77).} %68
 it is plausible to assume that it involved the
preparation of the bone or shell.\footnote[69]{The fact that T does not apear on oracle bones bearing the marginal notations ``sent in” (ju or lai) suggests a significant difference between those sent in and those brought in ( ( ). If T did involve preparation, it is possible that the scapulas or plastrons sent in were already prepared (cf. n. 52), and that those brought in were not, and had to be prepared by the \pinyinterm{shang} themselves.} %69
 Since this preparation was frequently carried out or
directed by one of the royal women,\footnote[70]{E.g., the inscriptions at \$436.3-4; 437.4. For preliminary studies of the royal women, see Hayashi Minao (1968), pp. 50-67; \pinyinterm{zhou-hongxiang} (1970/71); Matsumaru (1970), pp. 77–78, provides a good summary. It is not certain that these persons with the title Fu mou were all consorts of the king; some may have been the wives of the to-tzu +, ``the many sons'' or ``the many Tzu'' of the royal Tzu lineage (Shirakawa [1954], pp. 26-27). 16 ttant VE}%70
 and since the graph (though not necessarily the word)
% source: scan 034, printed 17
乙丁五三琤?
6st
NH+

NIKK 6+
was the same as the graph for ``altar stand,'' it is possible that the preparation was as much
ritual as mechanical. As a tentative, functional translation, therefore, I shall use the phrase
``ritually prepare.\footnote[71]{Cf. \pinyinterm{jiabian} 2071, \pinyinterm{jiabian} k'ao-shih, pp. 259-260. In contexts of ancestral sacrifice (\$150.1151.2), the graph referred to some form of altar stand on which to place offerings, and thus, by extension, to the ancestral spirits who descended to visit, or who were symbolized by, that altar (李 成-fu [1965], pp. 50-51; Ojima [1968], pp. 28-29; Katō [1970], 1155; cf. Tian Qianjun [1973]). The marginal notation may indicate, therefore, that the bone or shell was ``altared,'' that is, placed on such a stand, presumably to sanctify or consecrate it for its divinatory role. But note that although forms (a) T or (b) † could be used interchangeably for “altar stand,” form 6 was never used in the marginal notations (S151.3-152.3); this epigraphic rigor suggests that form a may not have meant ``altar stand'' when it was used in the marginal notations (on the possibility that identical graphs might have radically different meanings, see sec. 3.5.2). The possibility of another meaning besides ``sanctify,'' therefore, cannot be excluded. \pinyinterm{chen-mengjia-spaced} (1956), p. 18, suggested that in the marginal notations form a referred to the making of the hollows (cf. \pinyinterm{li-xueqin} [1959], p. 20); this is graphically plausible; the graph might depict a drill or gouger. I know of no later evidence to support such an interpretation, but \pinyinterm{shang} usage provides some support (see n. 74).}''
``71 Notations like ``Wo\footnote[72]{Wo could mean either “we” or the name of a statelet; frequently it is impossible to decide which meaning was intended. For an introduction to the problem, see \pinyinterm{chen-mengjia-spaced} (1956), pp. 314, 315; \pinyinterm{rao-short} (1959), pp. 705-706.}%72
 brought in 1,000 (shells); Fu 京 ritually
prepared 100. (Recorded by) Que'' and ``Wo brought in 1,000 (shells); Fu 京
ritually prepared 40. (Recorded by) Pin\footnote[73]{\pinyinterm{yibian} 6686 (S414.3); 6967 (S412.4); 7131 (S414.3). Cf. n. 42.}'' 73 indicate that though both plastrons were part
of a group of 1,000 which were probably obtained at one time, Fu 京 in the one case
prepared 100 of them, as noted by the divining official Que, and in the other case
prepared 40 of them, as recorded by Pin.\footnote[74]{Cf. Itō (1962a), p. 255. It seems unlikely that in contexts such as these T could mean ``gave as a gift, contributed''; if so, it would require that Wo conveyed the shells to Fu 京, who in turn ``gave'' part of them to the royal diviners. This seems improbable, for if the royal diviners received only 100 from Fu 京, why did they record that she was, as it were, holding back 900 more? The fact that (, the term for unpaired scapulas (n. 65), never follows ju, lai, or \&, suggests that single scapulas were not sent in or requisitioned. The term does follow the verb T, however, on every occasion that I have found it (e.g., \pinyinterm{yicun} 418; \pinyinterm{houbian} 2.33.10 [both S151.4]; Kikkō 1.18.14; \pinyinterm{cuibian} 1504 [both at S152.2, but Shima transcription may not be accurate for the second case]; \pinyinterm{yicun} 457), and this constitutes further evidence for thinking that I did not mean ``give'' or ``contribute,'' but may have meant ``prepare''; that is, scapulas would naturally be sent in pairs, two to an animal, but they could be prepared in even or odd numbers (cf. n. 71).}%74
 The number of shells involved in these two
cases was exceptionally large; the number of scapulas or plastrons prepared in this way was
usually 30 or less.\footnote[75]{\pinyinterm{renwen} shakubun, p. 353; Itō (1962a), p. 254.}%75

The officers who signed these notations sometimes have the same names as the court
diviners (sec. 2.2.1.1) of period I-成 (S110.2), Que (S414.3), Xuan (S315.1-2),
Pin (S276.2) and so on—and it is assumed that in these cases the diviners themselves were
responsible for recording the predivinatory notations.\footnote[76]{\pinyinterm{hu-houxuan} (1948), p. 317, provides a useful table of the officers' names. Preliminary research indicates that when the name was a diviner's it was generally not the diviner of that name who subsequently divined with the shell (e.g., \pinyinterm{pinbian} 267/268; 313/314; 377/378; 379/ 380; \pinyinterm{yibian} 6966/6967; 7126/7127), though the names might on occasion be the same (e.g., \pinyinterm{pinbian} 392/393). 17 ] 1 11 ] ( ] 1] ] ]}%76

KNKY \#13
(1
% source: scan 035, printed 18

\section[The Hollows]{\origsecnum{1.5}The Hollows}
\label{sec:chapters-ch01:1.5}

\subsection*{The Hollows}

After the marginal notations had been inscribed, a series of hollows was bored or
chiseled into the back or inner surface of the bone or shell.\footnote[77]{For evidence that the hollows were made after the marginal notations, see n. 60. Hollows were occasionally made on the front or outer surface of scapulas or plastrons (\pinyinterm{chen-mengjia-spaced} [1956], p. 11; \pinyinterm{zhang-bingquan} [1967], p. 854). For good examples, see \pinyinterm{jiabian} 2907; \pinyinterm{guo-moruo-spaced} (1972), pp. 8-11 (reproduced in Zhonghua renmin [1973], plates 35-36).}%77
 When heat was applied to the
hollow (sec. 1.6.2) in the back, the characteristic pu þ-shaped crack would form in front.
These hollows thinned the bone or shell, thus making it easier to crack; they determined
the crack's approximate shape and direction, and they determined its location. It is the
well-ordered hollows which distinguish the structured pyromancy of the \pinyinterm{shang} from the
free-form pyromancy of the Neolithic period.
The tools used to bore and chisel the hollows have not yet been specifically identified
in the \pinyinterm{anyang} finds. Presumably, the bored hollows were made with a bow drill of bronze
or stone, aided by the addition of an abrasive grinding powder. The chiseled hollows must
have been cut with some kind of jade knife, bronze chisel, or stone burin.\footnote[78]{What may have been a bronze ``drill,'' used to bore out the hollows, has been excavated at \pinyinterm{erligang}; its diameter and curvature match those of the circular hollows found in a scapula excavated at the same place (\pinyinterm{an-zhimin} [1954], p. 91, plate 16.1-2; cf. \pinyinterm{chen-mengjia-spaced} [1956], p. 12, plate 5; Kaizuka [1967], p. 202, fig. 206; \pinyinterm{liu-yuanlin} [1974], pp. 125–126). For a reconstruction of the \pinyinterm{shang} knife, based on the shape of the hollows, see \pinyinterm{zhang-bingquan} (1954), p. 219; on \pinyinterm{shang} engraving knives in general, see \ref{ch:2} (see ch. 2), n. 102, below. I can testify from personal experiment with modern steel blades and electric drills that making the hollows must have been difficult work. It was also exacting. The amount of shell left between the bottom of the oval hollow and the front surface of the plastron, for example, ranged between 0.2 and 0.5 mm. in thickness; a sharp implement and considerable skill must have been employed not to break through the shell (\pinyinterm{zhang-bingquan} [1954], pp. 219-220). It is uncertain if the inscriptions contain any reference to hollow making. \pinyinterm{guo-moruo-spaced} (\pinyinterm{cuibian} 1534, k'ao-shih, pp. 750-751) suggested that the graph * (S363.4) referred to the preparation of the bones, but this is speculative (cf. \pinyinterm{renwen} 1096, shakubun, p. 358, n.). The possibility that the T of the marginal notations referred to the making of the hollows has already been mentioned (nn. 71 and 74). Chan Pingleung (1972), pp. 39-41, elaborates the theory of Lao Kan (1957) that the upper part (4) of the graph for shih (oracle-bone form * ) represents the tool used to bore the hollows. This and his other interesting speculations are unsupported by the \pinyinterm{shang} inscriptions. The term used in \pinyinterm{zhou-dynasty} texts to refer to the boring or chiseling process was apparently ch'i (e.g., Hopkins [1919], pp. 115, 116-118; GSR 279b).}%78

The form of the hollows, whose shapes evolved with time, varied in accordance with
the thickness of the bone and shell. In general, the following rationale applied. When
the bone or shell was very thin, it might still be cracked in the Neolithic style by the application of heat to the unprepared surface.\footnote[79]{\pinyinterm{zhang-bingquan} (1967), p. 855; e.g., \pinyinterm{renwen} 978. Such cases are extremely rare in the \pinyinterm{xiaotun} finds.} %79
 When the bone or shell was thicker, single round
or oval hollows were used; single ovals frequently show no traces of having been burnt,
which suggests that some may have been uncompleted, and thus unused, double hollows.\footnote[80]{\pinyinterm{zhang-bingquan} (1954), p. 220. See too,\ref{sec:chapters-ch01:1.6.4} (see  sec. 1.6.4).} %80

When the material was still thicker, double hollows (in modern terminology, shuang-lien
yao-hsieh, ``double-linked pits,'' or tsuan-yü-tso, ``bored-and-chiseled'')
% source: scan 036, printed 19
were used (fig. 4). These were composed of \listitem{1} an oval or oblong-shaped hollow (what is
now called the tsao-ho, “date pit'' form) with the long axis vertical (that is, parallel
to the long dimension of the bone or shell); this was chiseled first and usually had a vertical
line at its deepest point, where the two sloping sides met; and \listitem{2} a slightly shallower,
semicircular or thumbnail-shaped hollow, either chiseled or bored, to the right or left of
the longer oval hollow. Generally, the axis of the chiseled oval was longer than the diameter
of the semicircular hollow.\footnote[81]{\pinyinterm{dong-short} (1929), p. 96; \pinyinterm{zhang-bingquan} (1954), p. 220, fig. 2; \pinyinterm{renwen} shakubun, introduction, p. 11; \pinyinterm{pinbian} 95, k'ao-shih, p. 128.}%81

\pinyinterm{xu-jinxiong} has established five categories for the bored-and-chiseled hollows
found on scapulas: \listitem{1} the round chiseled hollow bigger than, and containing, the long
chiseled hollow; \listitem{2} the small round drilled hollow; \listitem{3} the round chiseled hollow on the
shoulder of the long chiseled hollow; \listitem{4} the chiseled hollow located on the lower section of
the front of the bone; \listitem{5} the long chiseled hollow unaccompanied by a round chiseled
hollow.\footnote[82]{\pinyinterm{xu-jinxiong} (1974), p. 12; cf. 徐 (1973a), pp. 6, 40-41, 64-65. On the first two kinds of round hollows, see too \pinyinterm{liu-yuanlin} (1974), p. 129.} %82
 The way these different hollow shapes may be used as a criterion for periodic
classification is considered in\ref{sec:chapters-ch04:4.3.2.3} (see  sec. 4.3.2.3). With regard to the double hollows found on
scapulas, 徐 stresses that many of the so-called round hollows that previous scholars
have identified as drilled were, in fact, either chiseled or, more frequently, were simply
the result of heat chipping with no prior cutting; rubbings cannot reveal this distinction,
which must be made from an examination of the bone or shell itself. On shells, by contrast,
most of the long chiseled hollows were accompanied by round hollows that had been
chiseled.\footnote[83]{\pinyinterm{xu-jinxiong} (1974), pp. 6, 10-11, 13, 42, 65; cf. 徐 (1973a), pp. 15, 40-41, 88-89; on the distinction between bored and chiseled hollows, see too \pinyinterm{liu-yuanlin} (1974), pp. 125–127. n 0 ] 0 0}%83


\subsection[Hollow Placement]{\origsecnum{1.5.1}Hollow Placement}
\label{sec:chapters-ch01:1.5.1}

n
Cracks on shell were generally made with rigorous attention to their mirror
symmetry about the central axis, a symmetry which depended upon the intentional
regularity with which the hollows had originally been bored and chiseled. Since the way
the hollows and cracks were likely to be placed is one of the clues by which modern scholars
attempt to reconstruct whole scapulas or plastrons out of fragments (sec. 5.7), it is necessary
to understand \pinyinterm{shang} practice. The figures will make clear what is difficult to describe
simply or gracefully in words.
With few, if any, exceptions, the hollows on the rear of plastrons were cut so that the
round hollow was chiseled on the inside shoulder of the long chiseled hollow, that is, on
the shoulder closest to the central axis (fig. 4). As a result, the transverse cracks—which
formed in the round hollow (\ref{sec:chapters-ch01:1.6.2} (see see sec. 1.6.2))—all ran toward the central axis of the shell,
% source: scan 037, printed 20
those on the right side of the shell running to the left and vice versa.\footnote[84]{\pinyinterm{zhang-bingquan} (1956), pp. 240-241; \pinyinterm{dong-short} (1958), pp. 916-917, fig. 4. For examples, see \pinyinterm{pinbian} 1; 17; 23; 27; 33; 37, etc. Perfect symmetry was not always observed on the entoplastron; an odd number of hollows (usually one or three) might be carved into this bony plate, and in such cases the ``leftover'' hollow usually straddled or lay near the central axis of the shell (e.g., \pinyinterm{pinbian} 40; 42; 374; 378; 579); the entoplastron might also, in other cases, bear two hollows (e.g., \pinyinterm{pinbian} 46; 48) or none (it seems that when relatively few hollows were to be made on the shell, none at all were made on the entoplastron, e.g., \pinyinterm{pinbian} 355; 369 and \pinyinterm{yibian} 754; 4541, with ten hollows on each shell but none on the entoplastron). For a rare case in which the hollows on shell were chiseled asymmetrically, see 焦 (1961), p. 957, discussing \pinyinterm{yibian} 6692.}\footnote[85]{\pinyinterm{dong-short} (1949a), p. 266. E.g., \pinyinterm{pinbian} 65/66; \pinyinterm{yibian} 4679/4680.}\footnote[86]{Gibson (1938), p. 13, proposed that a bone-measuring instrument had been used to locate the site of the burning accurately, but no such instrument has been found. It has also been suggested that horizontal knife cuts found at the lower right edge of all the hollows on \pinyinterm{jiabian} 747 may have been initial guide marks for the hollow maker (\pinyinterm{liu-yuanlin} [1974], p. 117, plate 12.10).}\footnote[87]{On the meaning of right and left, see n. 48.}\footnote[88]{On the half-pair, see n. 65.}\footnote[89]{For examples of the various patterns, see: \pinyinterm{dong-short} (1936a), fig. 10 (3.3.0105), fig. 17 (\pinyinterm{jiabian} 657), fig. 19 (\pinyinterm{jiabian} 626); \pinyinterm{xu-jinxiong} (1973a), fig. 57b (\pinyinterm{kufang} 972 +976 [both drawings]), fig. 65 (\pinyinterm{kufang} 1009+1128+1132+ 1133+1153 [all drawings]), fig. 149 (\pinyinterm{jiabian} 2249).}
% source: scan 038, printed 21
had over a hundred double hollows chiseled into their backs,\footnote[90]{\pinyinterm{pinbian} 135 (fig. 4), for example, a particularly fine instance of precise, symmetrical placement, would have contained in its complete form 112 hollows chiseled into a plastron only 28.5 cm. long. \pinyinterm{pinbian} 184, 44 cm. long-the largest inscribed plastron found, and identified as Testudo emys (see appendix 1, sec. 3)—would have contained 204 hollows (董 [1948a], p. 122; \pinyinterm{chen-mengjia} [1956], p. 8).} %90
 but others of comparable
size had relatively few.\footnote[91]{\pinyinterm{pinbian} 9 (fig. 9), for example, had only ten hollows on a plastron 19.5 cm. long.}%91
 Small scapulas or plastrons might have had many or few hollows.
There were standard patterns (sec. 1.5.1) but there was no single standard number of
hollows per shell. Presumably, a relatively large plastron with few hollows indicates a time
when turtle shells were plentiful and could be prepared for wasteful use; a slightly smaller
plastron with many hollows indicates a time when shell was scarce and had to be prepared
for efficient use.\footnote[92]{\pinyinterm{pinbian} 4 and 6 are period I plastrons, both identified as Ocadia sinensis, both of approximately the same size (PL 27.5 and 29.5 cm., respectively), both recording divinations of the same diviner, and both excavated from the same pit; but the first, and slightly smaller, contains seventy-eight hollows; the second only twentyfour. Similarly, \pinyinterm{pinbian} 579 and \pinyinterm{yibian} 754 are both period I plastrons, both identified as Chinemys reevesi, both excavated from the same pit, and both bearing the same marginal notation, “Que sent in 250 (shells)'' (appendix 3, n. 11); the first is far smaller (PL 16.7 cm.) than the second (PL 22.4 cm.), yet the first has thirty-one hollows, the second only ten. This is not to argue that the hollow maker lacked certain numerical models; at least four of the shells from this shipment were of the identical ten-hollow pattern (\pinyinterm{pinbian} 355; 369; \pinyinterm{yibian} 754; 4541), at least nine were of the identical thirty-one-hollow pattern (\pinyinterm{pinbian} 357; 374; 389; 579; \pinyinterm{yibian} 3300; 7050; 7123; 7153; 7491); see too, the remarks about the entoplastron, n. 84.}%92


\section[Burning and Cracking]{\origsecnum{1.6}Burning and Cracking}
\label{sec:chapters-ch01:1.6}

At the moment of divination, the diviner proposed the charge and applied heat
to the hollow in the back; this scorched and charred the bone or shell with the result that
a pu-shaped crack formed on its front. The controlled application of heat sufficiently
fierce and enduring to crack the bone or shell systematically and presumably-given the
large number of cracks made per topic (sec. 2.5)-with some dispatch was no mean
technical feat.\footnote[93]{I have so far been unable to crack a scapula; presumably, I have not yet found the right hard wood or other heat source (or the requisite patience). I am not the only scholar who has had troubles of this sort. Takashima Kenichi describes an attempt to crack a fresh cow scapula in June 1969: Since we were ``serving barbecued beef at the party and had plenty of charcoal, I used a small piece of it to apply to a (which I made with an electric drill). ... Even after a few minutes, nothing happened. I thought then that a charcoal briquet didn't have enough heat; so I applied a red hot soldering iron. . . again nothing happened except for it getting a bit scorched. I got rather disgusted and so I just threw the damn thing in the whole mess of burning charcoal. . . . Divination was inauspicious.... In the meantime, while everyone was discussing why it wouldn't crack, the bone started to emit a stench. I was going to remove it from the fire and just before I could do so, the bone began to crack successively-Pak! Pak! Pak! It was terrific! Since we were also interested in the reconstruction of Archaic Chinese, we talked more about the sound it made: we had truly 'reconstructed' the Archaic Chinese * pak. . Karlgren, for one, wasn't too far off.... The lesson we then learned was that the prepared bones, and presumably shells, must have been dried pretty well'' (letter of 14 March 1971). Takashima, in short, reproduced primitive, un21 ] ] ] 10 0 D ]] ] 1 ] ] ]}%93

% source: scan 039, printed 22

\subsection[The Heat Source]{\origsecnum{1.6.1}The Heat Source}
\label{sec:chapters-ch01:1.6.1}

The carbonization marks on the hollow shoulders or in the hollows indicate that
the \pinyinterm{shang} diviners cracked the bone or shell with a discrete heat source. The heat source
was usually round and must have had approximately the same diameter as the hollow, for
the entire bottom of the hollow is usually scorched;\footnote[94]{\pinyinterm{zhang-bingquan} (1954), p. 222 and fig. 2C. \pinyinterm{xu-jinxiong} (1974), pp. 210-211, has noted that the burning sometimes left two circles on the bone. ``The inside circle is the area actually singed and the outside circle was affected by the heat but not directly singed'' (for examples, see the backs of Menzies 1857; 2639). For an unusual instance of small, oblong-shaped burn marks, indicating a heat source oblong in cross-section, see the back of the scapula fragment published as \pinyinterm{bali} 10 (D). For the possibility that the changing nature of the burn marks themselves may be used as a supplementary dating criterion, see sec. 4.3.2.5.} %94
 further, the heat must have been
intense, for it frequently destroyed the hollow's edge in whole or in part.\footnote[95]{徐 (1973a), pp. 15, 88-89; cf. sec. 4.3.2.5. There is no support for the view that the \pinyinterm{shang} diviner put the bone or shell directly in the fire, as argued, for example, by Ju (1969), p. 81; such a view has long been untenable (see, e.g., Chavannes [1911], p. 132).}''
The exact nature of the heat source has not yet been established. On the basis of
\pinyinterm{zhouli} passages, not notable for their clarity, it has been suggested that the tip of a brightly
glowing hardwood brand, kept on the verge of flaming by breaths of air, was applied to
the hollow.\footnote[96]{CLCY, ch. 47, p. 14b; ch. 48, p. 4b (Biot [1969], vol. 2, pp. 75, 78). For discussion, see Chavannes (1911), p. 132; \pinyinterm{dong-short} (1929), 102; \pinyinterm{zhang-bingquan} (1967), p. 858.}%96
 Various late \pinyinterm{zhou-dynasty} or \pinyinterm{han-dynasty} texts indicate that a thorn brand was used.\footnote[97]{For the texts, principally \pinyinterm{baihutong}, Shihchi, and \pinyinterm{yili}, see \pinyinterm{chen-mengjia-spaced} (1956), p. 12; \pinyinterm{zhang-bingquan} (1967), p. 858; cf. \pinyinterm{dong-short} (1929), pp. 102-103; Fujino (1960), pp. 13, 25. The fact that the sides of the hollows are not generally blackened indicates to 长 that charcoal was not used; Kaizuka (1967), p. 201, by contrast, believes that the \pinyinterm{zhouli} account does refer to something resembling charcoal. For the use of charcoal in the scapulimancy of other cultures, see Jochelson (1905), pp. 73-74; \pinyinterm{ling-chunsheng} (1934), p. 136. 胡 徐 (1782-1787), \ref{ch:4} (see ch. 4), pp. 3a-4b, reported that rural plastromancers in \pinyintext{jiangsu} used a paste composed of charcoal, white lead, and jujube to burn the shell (cf. 董 [1929], pp. 103-105); but the cracks reported by 胡 徐 are quite different from those found on \pinyinterm{shang} plastrons, and the burn marks on the oracle bones show the \pinyinterm{shang} did not use this method. Similar objections apply to the proposal of Matsumaru ([1959], p. 3) that the \pinyinterm{shang} placed some kind of moxa on the end of a bronze drill. According to Ban Nobutomo (17731846), hahaka wood-variously identified as that of the cherry, birch, or silk tree-was used on the island of Tsushima, continuing a tradition first} %97

prepared, Neolithic scapulimancy. Zhang Guangyuan made what is probably the most serious
modern attempt to reproduce \pinyinterm{shang} plastromancy
on 21 November 1973. It is hoped that his experiments will be the subject of a subsequent monograph; here I give only the briefest description
taken from his letter of 28 August 1975. He
prepared the double hollows with a knife, each
double hollow requiring thirty minutes. He then
placed three burning sticks of incense in the
hollow; after twenty-five minutes he heard the
first sound of cracking and heard two more crack
sounds at the end of thirty-five minutes. When he
turned the plastron over, the cracks were still
not visible on the front, and they did not appear
till the incense had been burning for fifty minutes.
He noted that the oil in the shell around the hollows
tended to burn, suggesting that the \pinyinterm{shang} must
have dried their shells before use. He also noted
that the cracks lost their visibility after five days
(something noted by 胡 徐 [1782-1787], \ref{ch:4} (see ch. 4),
p. 4b), a fact that helps explain the \pinyinterm{shang} custom
of carving the cracks (sec. 2.10). \pinyinterm{dong-zuobin}'s
accounts ([1929], p. 58; [1958], p. 921) of his
own plastromantic experiments are not specific
about the heat source used. Like Takashima and
长, however, he too remarked on the sound
that the crack made, the so-called kuei-yü
``voice of the turtle'' (first referred to, I believe,
by 胡 徐, loc. cit.). For why it would be desirable
to know the speed with which the \pinyinterm{shang} diviners
made individual cracks, \ref{sec:chapters-ch03:3.7.6} (see see sec. 3.7.6).
% source: scan 040, printed 23
More recent accounts of scapulimancy refer to clumps of flaming tinder or moxa, but
these would have given heat far more dispersed than that apparently applied to the \pinyinterm{shang}
bones.\footnote[98]{A clump of flaming tinder was used by Yünnanese scapulimancers in 1962, working with unprepared sheep scapulas with no hollows: ``The 'diviner' knelt beside the fire pit, ordered people to bring a bowl, a small clump of huo-ts'ao k (literally, ``fire grass''), a pinch of rice, and a normally well preserved sheep scapula. He placed the scapula in the bowl, held the socket with his left hand, and with his right rubbed the clump of huo-ts'ao along the edge of the scapula, at the same time muttering a prayer.... After the prayer, he took the clump of huo-ts'ao and used spit to stick it to the front of the bone, afterwards bringing fire to light it.... In less than a minute, there was a sound of the bone being split by burning. The diviner quickly brushed away the ashes of the straw, took a little black ash from the edge of the fire, and dabbed some on the back of the bone where it had already cracked; he then used more spit to wash it off'' (Lin Sheng [1964], pp. 98-99. Lin makes a similar report [1963], p. 164: diviners in southwest China rolled artemesia leaves or tinder between the hands to form a splint, which was placed, flaming, against the bone surface). This is unprepared, Neolithic divination; it is not \pinyinterm{shang}.} %98
 In short, our only \pinyinterm{shang} evidence, the round burn marks in the hollows, suggests
that the diviners used an intensely hot, round, hard object.\footnote[99]{Cf. \pinyinterm{xu-jinxiong} (1974), p. 256, on the burn marks found in period IV hollows: ``It seems that a very hard, perhaps metal, object was pressed firmly onto the bone.'' It is not certain that a bronze poker, such as the bore found at \pinyinterm{erligang} (n. 78), was used to provide the heat. Neolithic scapulimancers had already found a way to apply concentrated heat to bored hollows before the development of metallurgy; if a nonmetallic heat source was effective then, it presumably continued to be used in \pinyinterm{shang} times, have though technical improvements may occurred.}%99


\subsection[The Application of the Heat]{\origsecnum{1.6.2}The Application of the Heat}
\label{sec:chapters-ch01:1.6.2}

In the double hollows, the heat source was generally applied to the floor of the
semicircular side hollow. When an oval hollow alone was being scorched, the heat was
applied to one of the sloping shoulders of the hollow,\footnote[100]{The heat source might occasionally burn both shoulders, but this was apparently accidental (\pinyinterm{xu-jinxiong} [1973], pp. 15, 89).} %100
 at times causing, as we have seen,
the adjacent bone or shell surface to flake away.\footnote[101]{\pinyinterm{dong-short} (1929), pp. 96, 112; \pinyinterm{zhang-bingquan} (1954), p. 222, fig. 2B and C; (1967), p. 854. \pinyinterm{xu-jinxiong} (1973a), pp. 19, 102, indicates that ``because more heat was required to crack the bone [when the long hollow was used alone], the singed area was larger and chipped away more readily'' than in the case of the double hollows (cf. 徐 [1974], p. 82).} %101
 The vertical crack appeared behind
the center of the long chiseled hollow. Since, in the case of a true double hollow, the horizontal or transverse crack appeared behind the semicircular side hollow,\footnote[102]{Generally, the transverse crack runs right through the center of the burn mark. Cracks visible on the front of the bone are frequently invisible in back (e.g., Menzies 2090; 2455; 2476; 2639); when they are visible and viewed from the back, the transverse crack is frequently more visible than the vertical (e.g., Menzies 1857; 2455 [top hollow]; 2664; 2692). Any study of cracks should depend on observation of the actual bones, for cracks visible in the bone may not appear in rubbings (e.g., Menzies 2455; 2528; 2664; cf. \ref{ch:5} (see ch. 5), n. 54, below). 23 \{] | J} 102 it was the
recorded in Kojiki (Ban [1907], pp. 448, 449, 451;
Fujino [1960], p. 12; Philippi [1968], pp. 52, 82–
83, 559. Kaizuka [1967], p. 201, and Shirakawa
[1972], p. 20, without explanation read the
word as hahaki, some kind of twig or grass broom
[?]). Shirakawa (loc. cit. and [1970], p. 657)
suggests that at the critical moment of burning,
water was sprinkled on the shell; presumably this
generated steam and created greater stress. The
idea is plausible, but is supported by no \pinyinterm{shang},
\pinyinterm{zhou-dynasty}, or later Chinese evidence. For other
Japanese traditions, including that of inserting
skewers of bamboo into deer scapulas, see Fujino
(1960), p. 15.
% source: scan 041, printed 24
drilling or chiseling of the side hollow to the right or left of the oval one that determined
whether the pu crack would run to right or left. Similarly, the application of the heat to
the right or left shoulder of the single chiseled hollow determined the general direction of
the transverse crack. 103 The absence of a semicircular chiseled hollow at the side may also
have caused the end of the transverse crack to fork. 104 In short, the position of the side
hollow relative to the long oval and, presumably, the angle of the glowing brand as it was
applied to either the single or double hollow affected the angle of the pu crack. But it is
far from certain that the rudimentary manipulation this allowed would have significantly
altered the diviner's readings.\footnote[105]{The question of how the cracks were read pien 6385; 7456; 7797.} 10

\subsection[The Order of Burning]{\origsecnum{1.6.3}The Order of Burning}
\label{sec:chapters-ch01:1.6.3}

The crack numbers (sec. 2.4) subsequently engraved beside the cracks presumably
reveal the order in which the hollows were burned. 106 Five major patterns may be
described: 107
I. From top to bottom. The hollows at the top of the scapula or plastron were
burned first, followed in sequence by those below. Small plastrons tended to have only
one such vertical row of descending hollows and cracks per right and left side; large
plastrons might have several. 108
2. From inside to outside and from top to bottom. This is the commonest pattern
found on the plastrons. The hollow closest to the spine was burned first, followed in sequence
by those lying to its outside; when one horizontal row of hollows had been used in this
way, the burning continued in the next row down, starting again from the innermost
hollow. 109 The hollows of the horizontal row were not always cracked at one time; that
is, instead of having four hollows running from the center to the edge, with corresponding
Io3. \pinyinterm{chen-mengjia} (1956), p. II.
104. \pinyinterm{xu-jinxiong} (1973a), pp. 19, 102:
``Many bones have subsidiary cracking near the
end of the diagonal crack rather than one clear
diagonal crack. In most cases this occurred when
the bone was singed directly on or beside the long
chiseled hollow in the absence of a round chiseled
hollow. As the bone then lacked a thin area in
which the diagonal crack could form easily,
subsidiary cracking took place'' (cf. 徐 [1974],
p. 82). Forked cracks may be found on scapula
fragments such as Menzies 2639 (based on direct
observation). The fact that they also appear on
plastrons with their double hollows (e.g., \pinyinterm{pinbian}
8.4; 28.4; 49.11 [cracks 8 and 9]; 76.5) suggests
that, at least in the case of shell, the absence of the
round chiseled hollows was not required for
bifurcated cracks to form.
will be taken up in Studies. Shape alone was not the
determining factor; see \pinyinterm{zhang-bingquan} (1954).
106. Since the order in which the inscriptions
were carved presumably approximated the order
of burning, this section should be read in conjunction with sec. 2.9.4.
107. \pinyinterm{zhang-bingquan} (1956), pp. 231–236;
(1967), pp. 861, 867. Li Daliang (1972), pp. 8–30,
identifies eleven patterns of inscription placement
and sequence on plastrons; these patterns, which
relate to the sequence in which the hollows were
burned, may be used to refine Zhang Zongdong's categories.
108. E.g., one row per side: \pinyinterm{pinbian} 69; 354;
\pinyinterm{yibian} 2683; several rows per side: \pinyinterm{pinbian} 373;
\pinyinterm{yibian} 5279. For scapulas which follow pattern
I see, for example, \pinyinterm{dong-short} (1945), pt. 2, ch. 6,
p. 7b; ch. 7, p. 2a.
109. E.g., \pinyinterm{pinbian} 61; 63; 280; 293; 351; 乙-
% source: scan 042, printed 25
crack numbers 1, 2, 3, 4, one might have 1 and 2 for one charge, then 1 and 2 for another
charge.110
3. From outside to inside and from top to bottom. The hollow closest to the edge
was burned first, followed in sequence by the hollows to its inside; the next row down was
burned in the same way. This pattern is found rarely on plastrons. 111
4. From bottom to top. This pattern is common on scapulas. 112 On plastrons it is
found most frequently on the bridge, but it might extend over much of the shell.113
5. No evident pattern. This is rare.
The pattern of burning on some plastrons might be the same for all the divinations
performed. Other plastrons, however, might contain a mixture of patterns. 115
Density of Use
The pattern of burning on some plastrons might be the same for all the divinations
performed. Other plastrons, however, might contain a mixture of patterns. 115

\subsection[Density of Use]{\origsecnum{1.6.4}Density of Use}
\label{sec:chapters-ch01:1.6.4}

The lack of uniform density with which the hollows were prepared has already
been noted (sec. 1.5.2). There was similar variation in the efficiency with which the hollows
were burned. This is true, at least, in period I, for which we have the large number of
plastrons (reproduced in \pinyinterm{pinbian} and \pinyinterm{yibian}) needed for such a study. In some cases,
the hollows were all, or nearly all, used, 116 but this was not always true. For example,
there are 174 hollows in \pinyinterm{pinbian} 184, the unique tortoise plastron; only 43 of these hollows
have been burned, the remaining 131 are unused. Why, after arduous preparation, only
25 percent of the hollows were used is hard to explain, especially since the front of the
plastron, with the exception of a strip of inscriptions across the femoral laminae, is a tabula
rasa, with ample room for further divination charges. 117 It should be noted that in this
case 45 of the 174 hollows are single long hollows, and that not one of these 45 single hollows
was burned. 118 Conceivably, these were unburned because the shell was thought too thick
for the heat to be able to crack the single hollow (cf.\ref{sec:chapters-ch01:1.5} (see  sec. 1.5)).119 If so, it is curious that the
110. E.g., \pinyinterm{pinbian} 304.22-24.
III. E.g., \pinyinterm{pinbian} 90, k'ao-shih, p. 122; \pinyinterm{pinbian} 264.
112. E.g., \pinyinterm{jiatu} 128; cf. sec. 2.9.4.
113. E.g., on bridge: \pinyinterm{yibian} 2596; extended:
\pinyinterm{yibian} 5399 (fig. 5).
114. E.g., \pinyinterm{pinbian} 184; \pinyinterm{yibian} 811; 2244.
115. \pinyinterm{yibian} 3343, for example, contains patterns 1 and 2. The sequence on \pinyinterm{yibian} 5399 was
basically from bottom to top, inside to outside, but
three of its hollows were burned from outside to
inside (fig. 5); this was probably a diviner's error.
For other examples of confused order, see Li Daliang (1972), pp. 27-30, cf. pp. 87-88.
116. E.g., \pinyinterm{pinbian} 355 (ten hollows, all
burned); 374 (thirty-one hollows, all burned);
(twenty-five [?] hollows, twenty-five burned);@@
\pinyinterm{yibian} 754 (ten hollows, all burned); these period
I plastrons all bear the marginal notation ``Que
sent in 250 (shells)'' (see appendix 3, n. 11 and
were presumably divined at about the same time.
117. \pinyinterm{dong-short} (1948a), p. 122, estimated that on
the reconstructed plastron there would have been
hollows, of which 50 would have been used@@
He also cited (p. 127) a carapace (\pinyinterm{yibian} 6667)
with 108 hollows, of which only 68 were used.
\pinyinterm{zhang-bingquan} (1954), p. 225, lists a large
number of unburned hollows on the fragments
excavated in the thirteenth dig at \pinyinterm{xiaotun}.
His plate 2 (facing p. 225) gives pictures of burned
and unburned hollows. Numerous other examples
may be found in the figures of \pinyinterm{xu-jinxiong}
(1973a), using his key at pp. 40-41\%; see too, the
figures in \pinyinterm{xu-jinxiong-tight} (1974).
118. \pinyinterm{pinbian} 184, k'ao-shih, p. 266.
119. For similar cases, see \pinyinterm{pinbian} 246, k'ao-shih, p. 321; \pinyinterm{liu-yuanlin} (1974), p. 117, on
\pinyinterm{jiabian} 726.
% source: scan 043, printed 26
hollow makers did not continue their work by converting the single hollows into double
ones. In some cases, ritual error may also have been involved, but a detailed study of why
some hollows were burned and others were not is still needed.\footnote[120]{\pinyinterm{pinbian} 58 provides a particularly in teresting example: thirty double-linked hollows were bored and chiseled, and eight of these were not burned; of these eight, six were arranged symmetrically, three on the right and three on the left. It is significant that in the case of two out of the remaining twenty-two hollows that were burned and produced cracks, the cracks so produced were not numbered (for crack numbers, \ref{sec:chapters-ch02:2.4} (see see sec. 2.4)); and it is precisely the two unused hollows on the other side of the median line which correspond to those two hollows whose cracks were numberless. Figure 6 will make this clear. As \pinyinterm{zhang-bingquan} (\pinyinterm{pinbian} 58, k'ao-shih, pp. 93-94) suggests, it seems that after the crackmaking had taken place in these two hollows on the left side, some irregularity was discovered; hence the divinations were rejected, the cracks were not numbered, and the two corresponding hollows on the right were not used at all. This suggests that \pinyinterm{shang} divination had its strict rules.} %120
 The plausible hypothesis
that the proportion of hollows actually burned became smaller as the number of hollows
and the size of the plastron increased does not explain the evidence satisfactorily.\footnote[121]{The hypothesis works in some cases but not in others. For examples of plastrons with over twenty hollows and few burned, see \pinyinterm{pinbian} 389 (PL 18.3 cm., thirty-one [?] hollows, four burned); 579 (PL 16.7 cm., thirty-one [?] hollows, thirteen burned); 586 (PL 17.0 cm., thirty-one [?] hollows, eight burned); \pinyinterm{yibian} 7049 (PL 17.4 cm., thirtyone hollows, ten burned); 7123 (PL 17.2 cm., thirty-one hollows, eight burned. But in some cases we find that on plastrons with over twenty hollows, many or all of the hollows could be burned: \pinyinterm{pinbian} 374 (PL 19.0 cm., thirty-one hollows, thirty-one burned); 539 (PL 18.0 cm. [?], twenty-five [?] hollows, twenty-five burned); 616 (PL 18.2 cm., thirty-three hollows, twentyfive burned); \pinyinterm{yibian} 7153 (PL 17.0 [?] cm., thirty-one [?] hollows, twenty-six [?] burned); all these plastrons bear the same notation, “Que sent in 250 (shells)'' (see appendix 3, n. 11). Similarly, one can find shells with small numbers of hollows, few of which are burned, e.g., \pinyinterm{pinbian} 359 (PL 16.0 [?] cm., thirteen hollows, two burned).} %121

Among the factors which would have to be considered in any study of hollow use are
\listitem{1} the limiting role exercised by hollow, crack, and inscription density; \listitem{2} the nature of
the divination topic; \listitem{3} the thickness of bone or shell, that is, the ease or difficulty with
which it could be cracked; \listitem{4} the supply of prepared bones and shells available: \listitem{5} the
success or failure of the particular bone or shell in previous divinations;\footnote[122]{John Cikoski has suggested in conversation that certain bones and shells may have been little used because they had proved to be inaccurate; a study of the crack notations and prognostications and verifications (if any), which appear on shells that had been used inefficiently might provide a way to test this hypothesis, but the evidence will frequently be lacking. \pinyinterm{pinbian} 184, for example (see n. 118), contains ten charges about sacrifice, accompanied by four auspicious crack notations. But, as with all sacrifice divinations, no prognostications (sec. 2.7) or verifications (sec. 2.8) record the result.} %122
 \listitem{6} the bureaucratic problem of filing; that is, it may have been simpler to use a fresh plastron rather than
search out an old, partly used one; the rules of hollow use may have been no more regular
than those by which a modern clerk uses scratch paper.
It seems clear, in any event, that hollows were cut according to one set of criteria and,
at least on occasion, were burned according to another. Such a system was economically
inefficient, the diviners making no attempt to take full advantage of the considerable labor
expended by the hollow makers. Its spiritual efficiency, involving matters of ritual and
belief, we cannot judge. But it is evident that the supply of bone, shell, and labor was
sufficiently ample during the reign of \pinyinterm{wu-ding} for the diviners to have enjoyed great
% source: scan 044, printed 27
freedom in the way they made use of the divining materials. Just as \pinyinterm{wu-ding} lavished
sacrificial wealth upon his ancestors, his diviners showed no parsimony in their pyromantic
attempts to foreknow the future and influence the wishes of the spiritual powers.
